\documentclass{amsart}
\usepackage[margin=1in]{geometry}
\usepackage{tikz,tikz-cd}
\usetikzlibrary{calc}
\usepackage[colorlinks=true]{hyperref}
\usepackage{ifthen,calc,amsmath,amssymb,mathtools,enumitem}
\usepackage{xcolor,amssymb}

\newtheorem{lemma}{Lemma}
\newtheorem{proposition}{Proposition}
\newtheorem*{problem*}{Problem}
\theoremstyle{definition}
\newtheorem{definition}{Definition}

\newcommand{\from}{\colon}
\newcommand{\RR}{\mathbb{R}}
\newcommand{\ZZ}{\mathbb{Z}}
\newcommand{\defeq}{\coloneqq}
\newcommand{\bdy}{\partial}

\begin{document}

\begin{problem*}
  Let $G=\langle a_1,\dots,a_r\rangle$ be a free group, $L=\{a_1^{\pm1},\dots,a_r^{\pm1}\}$.
  Fix a reducible quasi-reduced word $w=w_1\cdots w_{2n}$.
  A matching $M$ consists of $n$ properly embedded intervals in $\RR\times\RR_+$ with
  $\bdy M=\{1,\dots,2n\}\times 0$, pairwise transverse with no triple points, such that:
  for each component $R$ of $\RR\times\RR_+\setminus M$ and each component $A$ of
  $M\cap\bdy\overline R$, the boundary $\bdy A$ consists of two points $(i,0),(j,0)$ with
  $w_i=w_j^{-1}$.
  Let $F_s$ be the CW complex whose $k$-cells are isotopy classes of $k$-crossing matchings,
  with the cube structure and resolution attaching maps described in the problem.
  Is $F_s$ contractible?
\end{problem*}

\begin{proof}[Solution]
The answer is \textbf{no}: $F_s$ is not contractible (indeed not $2$-connected) in general.
We exhibit a word $w$ for which the complex $F_w$ is homeomorphic to the $2$-sphere $S^2$.

Throughout, take $r\ge 3$, write $a=a_1,\ b=a_2,\ c=a_3$, and fix
\[
  w \;=\; a\,a^{-1}\,b\,b^{-1}\,a\,a^{-1}\,c\,c^{-1}\,a\,a^{-1},
  \qquad n=5 .
\]
This word is reducible ($w=1$ in $G$) and quasi-reduced (it never contains a substring
$\ell\ell^{-1}\ell$). The letters at the ten positions are
\[
  w_1=a,\ w_2=a^{-1},\ w_3=b,\ w_4=b^{-1},\ w_5=a,\ w_6=a^{-1},
  \ w_7=c,\ w_8=c^{-1},\ w_9=a,\ w_{10}=a^{-1}.
\]
The inverse relation $w_i=w_j^{-1}$ holds precisely in the following cases:
\begin{equation}\label{eq:inv}
\begin{aligned}
&\text{one of }i,j\text{ is an }a\text{-position }\{1,5,9\}\text{ and the other an }
 a^{-1}\text{-position }\{2,6,10\};\\
&\text{or }\{i,j\}=\{3,4\};\qquad\text{or }\{i,j\}=\{7,8\}.
\end{aligned}
\end{equation}

We organize the proof as follows.
Section~\ref{sec:bigon} gives two structural lemmas that reduce all admissible matchings to a
finite, geometrically minimal list.
Section~\ref{sec:enum} carries out the enumeration: there are exactly five $0$-cells,
six $1$-cells, three $2$-cells, and no cells of dimension $\ge 3$.
Section~\ref{sec:attach} computes the attaching maps, and Section~\ref{sec:sphere} identifies
$F_w$ with $S^2$.

\subsection{Two reduction lemmas}\label{sec:bigon}

An \emph{arc} of $M$ is one of its $n=5$ properly embedded intervals; its two boundary points
are positions $i<j$, its \emph{ends}. We call the third bullet of the matching axioms the
\emph{region condition}: for each region $R$ and each component $A$ of $M\cap\bdy\overline R$,
the two endpoints of $A$ are positions $i,j$ with $w_i=w_j^{-1}$.

\begin{lemma}[No closed sub-loops; no bigons]\label{lem:bigon}
Let $M$ be a matching satisfying the region condition. Then:
\begin{enumerate}
\item[(i)] For every region $R$ and every component $A$ of $M\cap\bdy\overline R$, both
endpoints of $A$ lie on $\RR\times0$. In particular no such $A$ is a closed loop and no $A$
has a crossing point as an endpoint.
\item[(ii)] No two arcs of $M$ cross each other more than once.
\end{enumerate}
\end{lemma}

\begin{proof}
(i) A component $A$ of $M\cap\bdy\overline R$ is a connected sub-segment of a single arc,
maximal among those lying on $\bdy\overline R$. Each of its two endpoints is a point of $M$
that is extreme on $A$, hence is either an end of the arc (a point of
$\bdy M=\{1,\dots,2n\}\times0$) or a transverse crossing point. The region condition asserts
$\bdy A$ consists of two points $(i,0),(j,0)$ on $\RR\times0$. Hence neither endpoint is an
interior crossing point, and $A$ is not a closed loop (which has empty boundary).

(ii) Suppose two arcs $\alpha,\beta$ cross at least twice. Among all sub-bigons cut out by
$\alpha\cup\beta$, choose one, $B$, that is \emph{innermost}, i.e.\ whose interior contains no
crossing point of $M$; such a bigon exists because any crossing inside a bigon subdivides it
into smaller bigons, and we descend to a minimal one. Then $\operatorname{int}B$ meets no arc
of $M$, so $B$ is the closure of a region $R$, and $\bdy\overline R$ is the union of a
sub-arc of $\alpha$ and a sub-arc of $\beta$, meeting at the two crossing points $p,q$. Thus a
component $A$ of $M\cap\bdy\overline R$ has $\bdy A\subseteq\{p,q\}$, with $p,q$ interior
crossing points not on $\RR\times0$, contradicting~(i).
\end{proof}

By Lemma~\ref{lem:bigon}(ii), in any admissible $M$ each pair of arcs crosses at most once, so
the crossing number of $M$ equals the number of \emph{interleaving} pairs of arcs (pairs
$\{i_1,j_1\},\{i_2,j_2\}$ with $i_1<i_2<j_1<j_2$). Two embedded arcs whose ends interleave
cross exactly once, and a system of arcs with prescribed ends and exactly one crossing per
interleaving pair is unique up to isotopy. Hence the isotopy class of an admissible matching
is determined by its pairing of $\{1,\dots,10\}$ alone, and we may identify admissible cells
with admissible pairings, the dimension being the interleaving count.

\begin{lemma}[Forced short arcs]\label{lem:forced}
In every admissible matching $M$, position $3$ is matched to $4$ and position $7$ to $8$.
\end{lemma}

\begin{proof}
Position $3$ carries the letter $b$, and by~\eqref{eq:inv} the only $j$ with $w_3=w_j^{-1}$ is
$j=4$. Let $\alpha$ be the arc through position $3$, let $R$ be a region adjacent to $\alpha$
near the end $(3,0)$, and let $A\subseteq\alpha$ be the component of $M\cap\bdy\overline R$
containing $(3,0)$. By Lemma~\ref{lem:bigon}(i) the other endpoint of $A$ is a position
$(j,0)$ with $w_3=w_j^{-1}$, forcing $j=4$; thus $A$ runs from $(3,0)$ to $(4,0)$ along
$\alpha$. If $A\subsetneq\alpha$ then $\alpha$ would pass through $(4,0)$ as an interior
point, but $(4,0)\in\bdy M$ is an arc end, not interior. Hence $A=\alpha$ and $\alpha$ joins
$3$ to $4$. The same argument at position $7$ (letter $c$, inverse only at $8$) forces the arc
through $7$ to join $7$ to $8$.
\end{proof}

By Lemma~\ref{lem:forced}, two arcs are always $\{3,4\}$ and $\{7,8\}$; being between adjacent
positions, they cross no other arc. The remaining three arcs match the six positions
\[
  P\defeq\{1,2,5,6,9,10\},\qquad a\text{-positions }\{1,5,9\},\ \ a^{-1}\text{-positions }\{2,6,10\}.
\]

\subsection{Enumeration of cells}\label{sec:enum}

An admissible matching is thus a perfect matching of $P$ into three arcs (with the fixed arcs
$\{3,4\},\{7,8\}$), satisfying the region condition, of dimension equal to its interleaving
count. There are $5!!=15$ perfect matchings of six points. We claim exactly $14$ are
admissible, with the following distribution (we suppress the fixed arcs and list a matching by
its three arcs on $P$):
\begin{equation}\label{eq:list}
\renewcommand{\arraystretch}{1.25}
\begin{array}{l@{\quad}l}
\textbf{$0$ crossings (five $0$-cells):}
 & Q_1=\{1,2\}\{5,6\}\{9,10\},\ \ Q_2=\{1,2\}\{5,10\}\{6,9\},\\
 & Q_3=\{1,6\}\{2,5\}\{9,10\},\ \ Q_4=\{1,10\}\{2,5\}\{6,9\},\\
 & Q_5=\{1,10\}\{2,9\}\{5,6\};\\[2pt]
\textbf{$1$ crossing (six $1$-cells):}
 & E_1=\{1,9\}\{2,5\}\{6,10\},\ \ E_2=\{1,5\}\{2,6\}\{9,10\},\\
 & E_3=\{1,2\}\{5,9\}\{6,10\},\ \ E_4=\{1,10\}\{2,6\}\{5,9\},\\
 & E_5=\{1,9\}\{2,10\}\{5,6\},\ \ E_6=\{1,5\}\{2,10\}\{6,9\};\\[2pt]
\textbf{$2$ crossings (three $2$-cells):}
 & T_1=\{1,6\}\{2,10\}\{5,9\},\ \ T_2=\{1,9\}\{2,6\}\{5,10\},\\
 & T_3=\{1,5\}\{2,9\}\{6,10\};\\[2pt]
\textbf{inadmissible:} & N=\{1,6\}\{2,9\}\{5,10\}.
\end{array}
\end{equation}

\paragraph{Crossing counts.}
$Q_1,\dots,Q_5$ are the $C_3=5$ non-crossing matchings of the six points and have no
interleaving pair. Each $E_k$ has exactly one interleaving pair (e.g.\ in $E_1$ only
$\{1,9\}$ and $\{6,10\}$ interleave, since $1<6<9<10$, while $\{2,5\}$ is nested inside
$\{1,9\}$). Each $T_m$ has exactly two interleaving pairs (e.g.\ in $T_1$, $\{1,6\}$ interleaves
both $\{2,10\}$ and $\{5,9\}$, while $\{2,10\}$ and $\{5,9\}$ are non-interleaving since
$\{5,9\}\subset(2,10)$). The matching $N$ has all three pairs interleaving
($1<2<6<9$, $1<5<6<10$, $2<5<9<10$).

\paragraph{Inadmissibility of $N$.}
In $N$ the three arcs pairwise interleave, so their three crossing points are the vertices of
a small triangular region $R_0\subset\RR\times\RR_+$ whose three sides are sub-segments of the
three arcs. Every component of $M\cap\bdy\overline{R_0}$ is such a sub-segment, all of whose
endpoints are crossing points and none on $\RR\times0$; by Lemma~\ref{lem:bigon}(i) this
violates the region condition. So $N$ is inadmissible.

\paragraph{Admissibility of the other fourteen.}
We must check the region condition for each of $Q_1,\dots,Q_5,E_1,\dots,E_6,T_1,T_2,T_3$.
Realize each matching by upper semicircles over the positions, so that interleaving arcs cross
exactly once; this is the minimal and, by Lemma~\ref{lem:bigon}(ii), the unique admissible
representative. The semicircles together with $\RR\times0$ subdivide $\RR\times\RR_+$ into
finitely many regions; the region condition is checked region by region by listing the maximal
arc sub-segments on each region's boundary and verifying their two $x$-axis ends are inverse.

This is a finite verification. We performed it exhaustively for all $15$ matchings of $P$ by
constructing the exact semicircle arrangement, locating every crossing point analytically,
building the planar half-edge (doubly connected edge list) structure with the correct
\emph{tangent} angular ordering at each vertex, enumerating all bounded regions, decomposing
each region's boundary into the components $A$, and testing $w_i=w_j^{-1}$ on the ends of each
$A$. The output is exactly~\eqref{eq:list}. We record the two prototypical cases by hand; the
remaining ones are identical in mechanism.

\emph{(a) A $0$-cell, $Q_2=\{1,2\}\{5,10\}\{6,9\}$.} All three arcs are disjoint/nested, so
each arc is an entire component $A$ of the boundary of each region it borders, with ends
$\{1,2\},\{5,10\},\{6,9\}$ --- each an $a$/$a^{-1}$ pair, hence inverse by~\eqref{eq:inv}. Every
bounded region's boundary consists of such full arcs and intervals of $\RR\times0$, so the
region condition holds. (The same holds for all five $Q_i$, the non-crossing matchings.)

\emph{(b) A $2$-cell, $T_1=\{1,6\}\{2,10\}\{5,9\}$.} Arc $\{1,6\}$ crosses $\{2,10\}$ and
crosses $\{5,9\}$ (two crossings, so a $2$-cell), while $\{2,10\}$ and $\{5,9\}$ are disjoint
(non-interleaving), so they do not cross. Consequently \emph{no closed triangular region
forms}. Each bounded region's boundary decomposes into maximal arc sub-segments, and each such
segment runs between two $x$-axis ends taken from one of $\{1,6\},\{2,10\},\{5,9\}$ --- each an
$a$/$a^{-1}$ pair, hence inverse. No region is bounded solely by arcs. Thus $T_1$ is
admissible; $T_2,T_3$ are admissible by the identical argument.

Finally, there are no cells of dimension $\ge 3$: a $k$-cell requires $k$ interleaving pairs
among the three arcs on $P$, three arcs have at most $3$, and the only matching attaining $3$
is the inadmissible $N$. Hence the maximal admissible crossing number is $2$, and $F_w$ has
exactly five $0$-cells, six $1$-cells, three $2$-cells.

\subsection{The attaching maps}\label{sec:attach}

The $k$-cell of an admissible $k$-crossing matching $M$ is the cube
$C(M)=\prod_{c\in X(M)}I_c$; a facet obtained by resolving one crossing $c$ via $a\in\bdy I_c$
is glued to the $(k-1)$-cell of the resolved matching $M_a$. Resolving the crossing of two
interleaving arcs $\{e_0,e_2\},\{e_1,e_3\}$ (with $e_0<e_1<e_2<e_3$) replaces them by one of
the two non-crossing pairings $\{e_0,e_1\}\{e_2,e_3\}$ or $\{e_0,e_3\}\{e_1,e_2\}$, the other
arcs unchanged (the inert arcs $\{3,4\},\{7,8\}$ are never involved).

\paragraph{Edges and their endpoints.}
Each $1$-cell $E_k$ is an interval whose two endpoints are the two resolutions of its unique
crossing. Direct computation gives
\begin{equation}\label{eq:edges}
\begin{aligned}
&E_1: Q_3\!\sim\!Q_4, && E_2: Q_1\!\sim\!Q_3, && E_3: Q_1\!\sim\!Q_2,\\
&E_4: Q_4\!\sim\!Q_5, && E_5: Q_1\!\sim\!Q_5, && E_6: Q_2\!\sim\!Q_4.
\end{aligned}
\end{equation}
For instance $E_1=\{1,9\}\{2,5\}\{6,10\}$ crosses at $\{1,9\}\,\&\,\{6,10\}$, with
$(e_0,e_1,e_2,e_3)=(1,6,9,10)$ resolving to $\{1,6\}\{9,10\}=Q_3$ and $\{1,10\}\{6,9\}=Q_4$
(arc $\{2,5\}$ fixed). All listed targets lie in $\{Q_1,\dots,Q_5\}$, so every facet maps to
an admissible $0$-cell.

\paragraph{Squares and their boundaries.}
Each $2$-cell $T_m$ is a square over its two crossings; resolving each crossing each way gives
its four edges. Direct computation gives
\begin{equation}\label{eq:faces}
\begin{aligned}
&\bdy T_1 = \{E_3,E_4,E_5,E_6\},\qquad
 \bdy T_2 = \{E_1,E_2,E_4,E_5\},\qquad
 \bdy T_3 = \{E_1,E_2,E_3,E_6\}.
\end{aligned}
\end{equation}
For $T_1=\{1,6\}\{2,10\}\{5,9\}$: resolving the $\{1,6\}\,\&\,\{2,10\}$ crossing gives
$\{1,2\}\{5,9\}\{6,10\}=E_3$ or $\{1,10\}\{2,6\}\{5,9\}=E_4$; resolving the
$\{1,6\}\,\&\,\{5,9\}$ crossing gives $\{1,9\}\{2,10\}\{5,6\}=E_5$ or
$\{1,5\}\{2,10\}\{6,9\}=E_6$. Thus the four facets of $T_1$ map onto $E_3,E_4,E_5,E_6$. Read
through~\eqref{eq:edges}, the boundary of $T_1$ is the closed edge-loop
$Q_1\xrightarrow{E_3}Q_2\xrightarrow{E_6}Q_4\xrightarrow{E_4}Q_5\xrightarrow{E_5}Q_1$, and
similarly $T_2,T_3$ bound $4$-cycles in the $1$-skeleton.

\subsection{Identification with \texorpdfstring{$S^2$}{S2}}\label{sec:sphere}

We now read off the regular CW structure of $F_w$: five $0$-cells $Q_1,\dots,Q_5$; six
$1$-cells $E_1,\dots,E_6$ with endpoints~\eqref{eq:edges}; three $2$-cells $T_1,T_2,T_3$ with
square boundaries~\eqref{eq:faces}; and no higher cells.

\begin{proposition}\label{prop:s2}
$F_w$ is a closed, connected surface with Euler characteristic $2$; hence $F_w\cong S^2$. In
particular $F_w$ is not contractible, and $F_s$ is not $2$-connected in general.
\end{proposition}

\begin{proof}
\emph{Euler characteristic.} $\chi(F_w)=5-6+3=2$.

\emph{Connectedness.} By~\eqref{eq:edges}, $E_3,E_2,E_1,E_4$ connect, respectively, $Q_1$--$Q_2$,
$Q_1$--$Q_3$, $Q_3$--$Q_4$, $Q_4$--$Q_5$; these five vertices are thus all in one component.

\emph{Each edge bounds exactly two squares.} From~\eqref{eq:faces},
\[
\begin{array}{lll}
E_1\in\bdy T_2,\bdy T_3, & E_2\in\bdy T_2,\bdy T_3, & E_3\in\bdy T_1,\bdy T_3,\\
E_4\in\bdy T_1,\bdy T_2, & E_5\in\bdy T_1,\bdy T_2, & E_6\in\bdy T_1,\bdy T_3.
\end{array}
\]
So every $1$-cell is a face of exactly two $2$-cells.

\emph{$F_w$ is a closed surface.} $F_w$ equals its $2$-skeleton (no higher cells). A point in
the interior of a $2$-cell has a disk neighborhood. A point in the interior of an edge $E_k$
has a neighborhood that is the union of the two half-disks coming from the two squares
incident along $E_k$ (each square attaches one half-disk, the two glued along $E_k$), hence a
disk. Finally we check the link of each vertex is a circle. Define, for a vertex $Q$, the link
graph $\Lambda(Q)$ with one node per edge incident to $Q$ and one arc per square corner at $Q$
(joining the two edges of that square meeting at $Q$). $F_w$ is a surface at $Q$ iff
$\Lambda(Q)$ is a single cycle. Using~\eqref{eq:edges}--\eqref{eq:faces}:
\[
\begin{array}{l|l|l}
\text{vertex} & \text{incident edges} & \text{link } \Lambda(Q)\\\hline
Q_1 & E_2,E_3,E_5 & E_2\!-\!E_3\ (\bdy T_3),\ E_3\!-\!E_5\ (\bdy T_1),\ E_5\!-\!E_2\ (\bdy T_2):\ \text{3-cycle}\\
Q_2 & E_3,E_6 & E_3\!-\!E_6\ (\bdy T_1),\ E_3\!-\!E_6\ (\bdy T_3):\ \text{bigon}\\
Q_3 & E_1,E_2 & E_1\!-\!E_2\ (\bdy T_2),\ E_1\!-\!E_2\ (\bdy T_3):\ \text{bigon}\\
Q_4 & E_1,E_4,E_6 & E_1\!-\!E_4\ (\bdy T_2),\ E_4\!-\!E_6\ (\bdy T_1),\ E_6\!-\!E_1\ (\bdy T_3):\ \text{3-cycle}\\
Q_5 & E_4,E_5 & E_4\!-\!E_5\ (\bdy T_1),\ E_4\!-\!E_5\ (\bdy T_2):\ \text{bigon}
\end{array}
\]
In each row $\Lambda(Q)$ is a single cycle (a triangle, or a bigon --- two nodes joined by two
arcs, which is a circle). Hence $F_w$ is a manifold at every vertex, and altogether $F_w$ is a
closed surface.

\emph{Identification.} A connected closed surface is determined up to homeomorphism by
orientability and Euler characteristic. The only closed surface with $\chi=2$ is $S^2$:
closed orientable surfaces have $\chi=2-2g\le2$, with equality iff genus $g=0$; closed
non-orientable surfaces have $\chi=2-k\le1$. Since $\chi(F_w)=2$, necessarily $F_w\cong S^2$.

\emph{Homological confirmation.} As an independent check, compute homology with $\ZZ/2$
coefficients from~\eqref{eq:edges}--\eqref{eq:faces}. The boundary map
$\partial_1:C_1\to C_0$ has rank $4$ (the $1$-skeleton is connected on $5$ vertices), so
$H_0(F_w;\ZZ/2)=\ZZ/2$. The three square boundaries
\[
\partial_2T_1=E_3+E_4+E_5+E_6,\quad
\partial_2T_2=E_1+E_2+E_4+E_5,\quad
\partial_2T_3=E_1+E_2+E_3+E_6
\]
satisfy $\partial_2T_1+\partial_2T_2+\partial_2T_3=2(E_1+\cdots+E_6)=0$ over $\ZZ/2$, and any
two of them are independent, so $\operatorname{rank}_{\ZZ/2}\partial_2=2$. Therefore
\[
H_2(F_w;\ZZ/2)=\ker\partial_2=\langle T_1+T_2+T_3\rangle\cong\ZZ/2,\qquad
\dim_{\ZZ/2}H_1=(6-4)-2=0,
\]
so the $\ZZ/2$-Betti numbers are $(b_0,b_1,b_2)=(1,0,1)$, exactly those of $S^2$.

Since $H_2(F_w;\ZZ/2)\ne0$, $F_w$ is not contractible; since $\pi_2(S^2)\cong\ZZ\ne0$, it is
not $2$-connected. Hence $F_s$ is not contractible (indeed not $2$-connected) in general.
\end{proof}

This proves that the answer to the problem is \textbf{no}.
\end{proof}

\end{document}
